\documentclass[10pt,a4paper]{amsart}
\usepackage{amsmath,amssymb,amsthm,graphicx}
% Retained minimal preamble for the Stacks Project chapter more-algebra.tex (More on Algebra).
% The chapter file pinned-original-source.tex carries only \input{preamble} on its first line, so no
% preamble text lives inside the chapter itself.  The Stacks Project uses one single shared preamble for
% every chapter; the snapshot stages it identically as sources/stacks-chow/preamble.tex and
% sources/stacks-etale-cohomology/preamble.tex.  The block below is copied mechanically, line for line,
% from sources/stacks-chow/preamble.tex lines 14-18 (the Xy-pic setup) and 154-217 (the native theorem
% declarations together with the mathematical macros).  Nothing was added, removed or redefined; the
% wrapper class and the amsmath/amssymb/amsthm/graphicx base packages are supplied by the assembler.
% Chapter: more-algebra.tex -- More on Algebra.
% For commutative diagrams we use Xy-pic
\usepackage[all]{xy}

% We use 2cell for 2-commutative diagrams.
\xyoption{2cell}

% Theorem environments.
%
\theoremstyle{plain}
\newtheorem{theorem}[subsection]{Theorem}
\newtheorem{proposition}[subsection]{Proposition}
\newtheorem{lemma}[subsection]{Lemma}

\theoremstyle{definition}
\newtheorem{definition}[subsection]{Definition}
\newtheorem{example}[subsection]{Example}
\newtheorem{exercise}[subsection]{Exercise}
\newtheorem{situation}[subsection]{Situation}

\theoremstyle{remark}
\newtheorem{remark}[subsection]{Remark}
\newtheorem{remarks}[subsection]{Remarks}

\numberwithin{equation}{subsection}

% Macros
%
\def\lim{\mathop{\mathrm{lim}}\nolimits}
\def\colim{\mathop{\mathrm{colim}}\nolimits}
\def\Spec{\mathop{\mathrm{Spec}}}
\def\Hom{\mathop{\mathrm{Hom}}\nolimits}
\def\Ext{\mathop{\mathrm{Ext}}\nolimits}
\def\SheafHom{\mathop{\mathcal{H}\!\mathit{om}}\nolimits}
\def\SheafExt{\mathop{\mathcal{E}\!\mathit{xt}}\nolimits}
\def\Sch{\mathit{Sch}}
\def\Mor{\mathop{\mathrm{Mor}}\nolimits}
\def\Ob{\mathop{\mathrm{Ob}}\nolimits}
\def\Sh{\mathop{\mathit{Sh}}\nolimits}
\def\NL{\mathop{N\!L}\nolimits}
\def\CH{\mathop{\mathrm{CH}}\nolimits}
\def\proetale{{pro\text{-}\acute{e}tale}}
\def\etale{{\acute{e}tale}}
\def\QCoh{\mathit{QCoh}}
\def\Ker{\mathop{\mathrm{Ker}}}
\def\Im{\mathop{\mathrm{Im}}}
\def\Coker{\mathop{\mathrm{Coker}}}
\def\Coim{\mathop{\mathrm{Coim}}}

% Boxtimes
%
\DeclareMathSymbol{\boxtimes}{\mathbin}{AMSa}{"02}

%
% Macros for moduli stacks/spaces
%
\def\QCohstack{\mathcal{QC}\!\mathit{oh}}
\def\Cohstack{\mathcal{C}\!\mathit{oh}}
\def\Spacesstack{\mathcal{S}\!\mathit{paces}}
\def\Quotfunctor{\mathrm{Quot}}
\def\Hilbfunctor{\mathrm{Hilb}}
\def\Curvesstack{\mathcal{C}\!\mathit{urves}}
\def\Polarizedstack{\mathcal{P}\!\mathit{olarized}}
\def\Complexesstack{\mathcal{C}\!\mathit{omplexes}}
% \Pic is the operator that assigns to X its picard group, usage \Pic(X)
% \Picardstack_{X/B} denotes the Picard stack of X over B
% \Picardfunctor_{X/B} denotes the Picard functor of X over B
\def\Pic{\mathop{\mathrm{Pic}}\nolimits}
\def\Picardstack{\mathcal{P}\!\mathit{ic}}
\def\Picardfunctor{\mathrm{Pic}}
\def\Deformationcategory{\mathcal{D}\!\mathit{ef}}

% --- appended 2026-10-12 (Chief): the project own definitions of the reference/slogan/history
% display environments, copied mechanically from the shared Stacks preamble lines 8-12
% (\usepackage{verbatim} plus the three \newenvironment lines, which the chapter preamble keeps
% verbatim from the project preamble). Nothing added, removed or redefined: these environments
% carry no mathematical content and the project itself typesets them as comments. Without these
% three lines a counted statement whose own text uses them fails with Environment slogan undefined.
\usepackage{verbatim}
\newenvironment{reference}{\comment}{\endcomment}
\newenvironment{slogan}{\comment}{\endcomment}
\newenvironment{history}{\comment}{\endcomment}
\usepackage{cleveref}
\crefname{originalsection}{original section}{original sections}
\Crefname{originalsection}{Original section}{Original sections}
\makeatletter
\newlabel{derived-section-generators}{{09SI}{source}{}{sourcepointer-1}{}}
\newlabel{derived-section-generators@cref}{{[originalsection][0][]09SI}{[1][0][]source}{}{}{}}
\makeatother
\crefname{originallemma}{original lemma}{original lemmas}
\Crefname{originallemma}{Original lemma}{Original lemmas}
\makeatletter
\newlabel{derived-lemma-representable-homological}{{0149}{source}{}{sourcepointer-2}{}}
\newlabel{derived-lemma-representable-homological@cref}{{[originallemma][0][]0149}{[1][0][]source}{}{}{}}
\makeatother
\title{The ghost lemma: a composition of maps that are zero on cohomology between stages of the generated subcategory vanishes}
\author{The Stacks Project Authors}
\date{stacks-project:more-algebra.tex:local-snapshot-20261012; GNU FDL 1.2 or later}
\begin{document}
\maketitle
\setcounter{subsection}{1}
\noindent\textit{Editorial note.} All delivered mathematical text is the author's own, copied from the pinned Stacks Project chapter file more-algebra.tex; no statement or proof step was rewritten, reordered or removed. The counted statement is the labeled lemma environment at source lines 20663-20677 (canonical tag 0FXH) and the counted proof is the proof environment at source lines 20679-20708, which immediately follows it and is its own proof.\par\medskip
\noindent\textit{Editorial note.} The counted statement contains the author's reference environment recording the origin of the result; the retained preamble defines that environment as comment-type, so its body and the citation inside it are discarded when the document is typeset, which is why the card needs no bibliography.\par\medskip
\noindent\textit{Editorial note.} The two references are to the section introducing the notation of the stages and to the representability statement used in the induction; the kinds recorded are the type words appearing in the labels themselves. Both are navigation stubs with their roles recorded in the reasons.\par\medskip
\noindent\textit{Editorial note.} The preamble span runs to line 90, the current end of the retained preamble, which is exactly what makes the reference environment a comment; the excerpt serves only as preamble, never as a span coordinate.\par\medskip
\noindent\textit{Editorial note.} The source URL is the source-level Stacks Project URL, as required when more than one card is produced from one source; the per-claim tag is carried by claim\_key and by this note.\par\medskip
\noindent\textit{Allowed standard prerequisites (editorial).}\begin{itemize}
\item An object of the first stage of the subcategory generated by an object is a direct summand of a finite direct sum of shifts of that object, and maps out of a shift of the ring into a complex are computed by the cohomology of the complex. used at proof lines 20679-20686. Reduces the case of one stage to the vanishing of a map whose induced map on cohomology is zero.
\item An object of the n-th stage of the generated subcategory sits in a distinguished triangle whose first term is of the first stage and whose third term is of the previous stage. used at proof lines 20689-20694. Provides the induction step, after replacing the object by the object of the triangle.
\item A morphism whose composite with the map to the third term of a distinguished triangle vanishes factors through that third term, by the representability statement for the homological functor. invoked at proof lines 20701-20703. Factors the first map of the chain and reduces the claim to the induction hypothesis applied to the remaining maps.
\end{itemize}\medskip
\ifdefined\hypertarget\hypertarget{sourcepointer-1}{}\fi
\noindent\textit{Source reference.} Original section 09SI: sources/stacks-chow/tags.txt, source lines 12356--12356. The reference points to the section where the notation for the stages of the generated subcategory is introduced, cited at proof line 20669; it fixes notation and the induction over stages is inside the counted span.\par\medskip
\ifdefined\hypertarget\hypertarget{sourcepointer-2}{}\fi
\noindent\textit{Source reference.} Original lemma 0149: sources/stacks-chow/tags.txt, source lines 1464--1464. The reference is applied at proof lines 20701-20703 to factor the first map of the chain through the third term of the distinguished triangle; the induction that consumes the factorization is inside the counted span.\par\medskip
\begin{lemma}
\label{lemma-ghost-lemma}
\begin{reference}
\cite{Kelly}
\end{reference}
Let $R$ be a ring. Let $n \geq 1$. Let $K \in \langle R \rangle_n$ with
notation as in Derived Categories, Section \ref{derived-section-generators}.
Consider maps
$$
K \xrightarrow{f_1} K_1 \xrightarrow{f_2} K_2
\xrightarrow{f_3} \ldots \xrightarrow{f_n} K_n
$$
in $D(R)$. If $H^i(f_j) = 0$ for all $i, j$, then
$f_n \circ \ldots \circ f_1 = 0$.
\end{lemma}

\begin{proof}
If $n = 1$, then $K$ is a direct summand in $D(R)$ of a bounded complex
$P^\bullet$ whose terms are finite free $R$-modules and whose differentials
are zero. Thus it suffices to show any morphism $f : P^\bullet \to K_1$
in $D(R)$ with $H^i(f) = 0$ for all $i$ is zero. Since $P^\bullet$ is
a finite direct sum $P^\bullet = \bigoplus R[m_j]$ it suffices to show
any morphism $g : R[m] \to K_1$ with $H^{-m}(g) = 0$ in $D(R)$ is zero.
This follows from the fact that $\Hom_{D(R)}(R[-m], K) = H^m(K)$.

\medskip\noindent
For $n > 1$ we proceed by induction on $n$. Namely, we know that $K$
is a summand in $D(R)$ of an object $P$ which sits in a distinguished triangle
$$
P' \xrightarrow{i} P \xrightarrow{p} P'' \to P'[1]
$$
with $P' \in \langle R \rangle_1$ and $P'' \in \langle R \rangle_{n - 1}$.
As above we may replace $K$ by $P$ and assume that we have
$$
P \xrightarrow{f_1} K_1 \xrightarrow{f_2} K_2
\xrightarrow{f_3} \ldots \xrightarrow{f_n} K_n
$$
in $D(R)$ with $f_j$ zero on cohomology.
By the case $n = 1$ the composition $f_1 \circ i$ is zero.
Hence by Derived Categories, Lemma \ref{derived-lemma-representable-homological}
we can find a morphism $h : P'' \to K_1$ such that $f_1 = h \circ p$.
Observe that $f_2 \circ h$ is zero on cohomology.
Hence by induction we find that $f_n \circ \ldots \circ f_2 \circ h = 0$
which implies $f_n \circ \ldots \circ f_1 = 
f_n \circ \ldots \circ f_2 \circ h \circ p = 0$ as desired.
\end{proof}

\end{document}
